% Source: https://artofproblemsolving.com/community/c3h1259727_right_triangle_lengths
% Stable identity: aops-topic-1259727-q01
\begin{asy}
usepackage("amsmath");
size(4cm);
defaultpen(fontsize(10pt));
pair O=(0,0),A=(4,0),B=(0,3); pair H=(A*9+B*16)/25;
draw(O--A--B--cycle,linewidth(1)); draw(O--H,dashed);
draw((0.3,0)--(0.3,0.3)--(0,0.3));
label("$O$",O,SW);label("$A$",A,SE);label("$B$",B,N);label("$H$",H,NE);
label("\mbox{}",(min(currentpicture,user=true)+max(currentpicture,user=true))/2,black); // keeps the text after the figure black when LaTeX typesets the labels inline
\end{asy}

In a right triangle, the altitude to the hypotenuse gives three similar triangles, so
\[OA^2=AH\cdot AB,\qquad OB^2=BH\cdot AB,\qquad OH^2=AH\cdot HB.\]
Multiplying, $(OA\cdot OB\cdot OH)^2=(AB\cdot AH\cdot HB)^2=21^2$, so $OA\cdot OB\cdot OH=21$ for every such triangle. Hence $M=m=21$ and $M+m=\boxed{42}$.
